% Originally: content/week-03.tex, \section{Cauchy--Schwarz} (Corsin Week 3, Friday)

\section{Cauchy--Schwarz}
\label{sec:cauchy_schwarz}

\subsection{Recap: inner products and norms}
\label{sec:inner_products_recap}

% Transition: Claude Opus 5 (High)
Corsin opens his Friday session by recalling the two definitions, which this chapter has already
stated in \cref{sec:structured_spaces}, where the hierarchy of structured spaces was drawn and
where the first supplements needed them. Rather than repeat them, here is what to have in mind:
\begin{itemize}
  \item An \emph{inner product} (\cref{def:inner_product}) is \textbf{bilinear},
    \textbf{symmetric} and \textbf{definite}. It measures length \emph{and} angle.
  \item A \emph{norm} (\cref{def:norm}) is \textbf{definite}, \textbf{homogeneous} and
    satisfies the \textbf{triangle inequality}. It measures length only.
  \item Each induces the next, via $\|v\| := \sqrt{\langle v,v\rangle}$ and
    $d(x,y) := \|x-y\|$, and neither implication reverses
    (\cref{prop:induced_norm_and_metric}, \cref{rem:norm_not_from_inner_product}).
\end{itemize}
What this section adds is the \emph{geometry}: why $\langle x,y\rangle$ deserves to be read as an
angle at all, and the resulting second proof of Cauchy--Schwarz.

\subsection{Geometric meaning of the inner product}

In $\mathbb{R}^2$, the following identity holds:
\[
\langle x,y\rangle := x_1y_1 + x_2y_2 = \|x\|\|y\|\cos\theta
\]
where $\theta$ is the angle between $x, y \in \mathbb{R}^2$.

To see this, conveniently choose our basis vectors such that $x = (x_1, 0)$, $y = (y_1, y_2)$.

\begin{center}
\begin{tikzpicture}[>=Latex, scale=1.1]
  \begin{scope}[shift={(0,0)}]
    \draw[->] (-0.3,0) -- (2.2,0);
    \draw[->] (0,-0.3) -- (0,2.2);
    \draw[blue!80!black, thick, ->] (0,0) -- (1.5,0.6) node[right] {$x$};
    \draw[orange!90!black, thick, ->] (0,0) -- (0.8,1.8) node[above] {$y$};
    \draw[purple, thick] (0.5,0.2) arc (21.8:66.0:0.5);
    \node[purple, font=\footnotesize] at (0.7,0.45) {$\theta$};
  \end{scope}

  \draw[->, thick] (2.4,1.0).. controls (2.8,1.3).. (3.3,1.0) node[midway, above, red, font=\tiny] {\shortstack{change of basis\\(rotation)}};

  % Same lengths and same angle as the left panel -- a rotation preserves both.
  % Left:  |x| = 1.616 at 21.80 deg, |y| = 1.970 at 66.04 deg, so theta = 44.24 deg.
  % Right: x rotated onto the axis at (1.616, 0), y at 44.24 deg with |y| unchanged.
  \begin{scope}[shift={(3.8,0)}]
    \draw[->] (-0.3,0) -- (2.5,0);
    \draw[->] (0,-0.3) -- (0,2.2);
    \draw[blue!80!black, thick, ->] (0,0) -- (1.616,0) node[below] {$x$};
    \draw[orange!90!black, thick, ->] (0,0) -- (1.411,1.375) node[above right] {$y$};
    \draw[dotted, gray] (1.411,1.375) -- (1.411,0);
    \draw[purple, thick] (0.5,0) arc (0:44.24:0.5);
    \node[purple, font=\footnotesize] at (0.72,0.22) {$\theta$};
  \end{scope}
\end{tikzpicture}
\end{center}

Then
\[
\langle x,y\rangle = x_1y_1 = x_1\|y\|\cos\theta = \|x\|\|y\|\cos\theta.
\]

For more rigour, you will show in Linalg that the rotations in $\mathbb{R}^2$ are given by
matrices $R$ satisfying $\langle Rx, Ry\rangle = \langle x,y\rangle$.

In a general inner product space $V$, this becomes the \textbf{definition} of the angle $\theta$.
It then makes sense to define the \newterm{projection onto $v$}, $\pi_v$:
\[
\pi_v(u) := \frac{\langle u,v\rangle}{\|v\|^2}\,v = \|u\|\cos\theta\,\frac{v}{\|v\|}.
\]

\begin{center}
\begin{tikzpicture}[>=Latex, scale=1.2]
  \draw[->] (-0.3,0) -- (3.5,0);
  \draw[->] (0,-0.3) -- (0,2.5);

  \draw[orange!90!black, thick, ->] (0,0) -- (3.2,0.8) node[right] {$v$};
  \draw[blue!80!black, thick, ->] (0,0) -- (1.2,2.0) node[above left] {$u$};

  \coordinate (P) at (1.60, 0.40);
  \coordinate (U) at (1.2, 2.0);

  \draw[purple!80!black, thick, dashed] (U) -- (P);
  \fill[purple!80!black] (P) circle (2pt) node[below right, font=\small] {$\pi_v(u)$};

  \draw[green!60!black, thick] (0.5,0.125) arc (14:59:0.5);
  \node[green!60!black, font=\small] at (0.6,0.3) {$\theta$};
\end{tikzpicture}
\end{center}

In words, $\pi_v(u)$ is the part of $u$ which points along $v$. In particular,
\[
\|\pi_v(u)\| = \|u\|\,|\cos\theta| = \frac{|\langle u,v\rangle|}{\|v\|}.
\]
In this sense, Cauchy--Schwarz says something visually obvious. Reading the picture above, all
three of the following are the same statement:
\[
\frac{|\langle u,v\rangle|}{\|v\|} \leq \|u\| \iff \|\pi_v(u)\| \leq \|u\| \iff \|u\|\,|\cos\theta| \leq \|u\|.
\]

\begin{remark}[Hint for the proof of Cauchy--Schwarz]
\label{rem:proof_hint_cauchy_schwarz}
Use definiteness and expand the term
\[
\langle u - \pi_v(u),\ u - \pi_v(u)\rangle = \dots
\]
where $u - \pi_v(u)$ is the \textbf{perpendicular part} of $u$ to $v$.
\end{remark}

% Sonnet 5 (Medium)
\subsection{The Cauchy--Schwarz inequality}

% Corsin's own notes stop at the hint above (rem:proof_hint_cauchy_schwarz) and never state the
% inequality itself as a formal result, which is notable given that it names the section. The
% theorem below fills that in, carrying out exactly the computation he hints at.

\begin{theorem}[Cauchy--Schwarz inequality]
\label{thm:cauchy_schwarz}
Let $V$ be an inner product space (\cref{def:inner_product}) with induced norm
$\|\cdot\| := \sqrt{\langle\cdot,\cdot\rangle}$ (\cref{def:norm}). Then for all $u,v \in V$,
\[|\langle u,v\rangle| \leq \|u\|\,\|v\|,\]
with equality if and only if $u,v$ are linearly dependent.
\end{theorem}

\begin{proof}
If $v = 0$, both sides vanish and the claim holds trivially, so assume $v \neq 0$. Following
\Cref{rem:proof_hint_cauchy_schwarz}, definiteness (\cref{def:inner_product}) gives
\[
  0 \ \leq\ \langle u-\pi_v(u),\ u-\pi_v(u)\rangle
    \ =\ \langle u,u\rangle - 2\langle u,\pi_v(u)\rangle + \langle \pi_v(u),\pi_v(u)\rangle.
\]
Using $\pi_v(u) = \tfrac{\langle u,v\rangle}{\|v\|^2}v$ and bilinearity, the two right-hand terms
are
\begin{align*}
  \langle u,\pi_v(u)\rangle &= \frac{\langle u,v\rangle^2}{\|v\|^2}, \\[0.3ex]
  \langle \pi_v(u),\pi_v(u)\rangle
    &= \frac{\langle u,v\rangle^2}{\|v\|^4}\langle v,v\rangle
     = \frac{\langle u,v\rangle^2}{\|v\|^2}.
\end{align*}
Note that the two are equal, so substituting them leaves a single copy behind:
\[
  0 \ \leq\ \|u\|^2 - \frac{2\langle u,v\rangle^2}{\|v\|^2} + \frac{\langle u,v\rangle^2}{\|v\|^2}
    \ =\ \|u\|^2 - \frac{\langle u,v\rangle^2}{\|v\|^2}.
\]
This rearranges to $\langle u,v\rangle^2 \leq \|u\|^2\|v\|^2$, that is,
$|\langle u,v\rangle| \leq \|u\|\,\|v\|$. By definiteness, equality holds if and only if
$u - \pi_v(u) = 0$, which says precisely that $u$ is a scalar multiple of $v$.
\end{proof}

% Generator: Claude Opus 5 (High)
\begin{remark}[The same proof twice]
\label{rem:two_proofs_cauchy_schwarz}
This is the second proof of Cauchy--Schwarz in these notes;
\cref{prop:cauchy_schwarz_geometric_proof} gave another, following Diego Torres
Tejeda. The two look different at first. One minimises the quadratic
$f(\lambda) := \|v-\lambda w\|^2$ over $\lambda$, while the other expands
$\|u-\pi_v(u)\|^2 \geq 0$. But the minimiser found in the first is
$\lambda_0 = \tfrac{\langle v,w\rangle}{\|w\|^2}$, which is exactly the coefficient defining
$\pi_w(v)$. Consequently, the first proof \emph{derives} the projection by optimisation, and the
second \emph{postulates} it and checks that it works. Keeping both is worthwhile precisely because
the pair shows where the projection comes from. It is not a lucky guess, but the closest point of
the line to $v$.
\end{remark}


% --- exercises relocated here from the dissolved sheets ---

% Originally: content/exercise-sheets/week-03.tex, problem 3.7
% Source: exercises/Ex3_Analysis2_eng.pdf

\begin{exercise}[3.7 --- Cauchy--Schwarz]
\label{ex:3.7}
Let $V$ be a vector space over $\mathbb{R}$, let $\langle\cdot,\cdot\rangle$ be an inner product
on $V$, and let $\|\cdot\| : V \to \mathbb{R}$ be given by $\|v\| = \sqrt{\langle v,v\rangle}$.
Then the inequality
\begin{equation*}
|\langle v,w\rangle| \leq \|v\|\|w\| \tag{1}
\end{equation*}
holds for all $v, w \in V$. Furthermore, equality in (1) holds if and only if $v$ and $w$ are
linearly dependent.
\end{exercise}

\begin{remark}
Corsin devotes the whole second half of Friday (\cref{sec:cauchy_schwarz}) to this problem,
including the geometric reading and an explicit proof hint.
\end{remark}

% Originally: content/exercise-sheets/week-03.tex, problem 3.8
% Source: exercises/Ex3_Analysis2_eng.pdf

\begin{exercise}[3.8 --- Inner product and norm]
\label{ex:3.8}
Let $V$ be a vector space over $\mathbb{R}$, let $\langle -,-\rangle$ be an inner product on $V$.
The map defined by
$$\|\cdot\| : V \to \mathbb{R}, \qquad \|v\| = \sqrt{\langle v,v\rangle}$$
satisfies the triangular inequality and is a norm.
\end{exercise}

% Originally: content/exercise-sheets/week-03.tex, problem 3.9
% Source: exercises/Ex3_Analysis2_eng.pdf

\begin{exercise}[3.9 --- All norms are equivalent in $\mathbb{R}^n$]
\label{ex:3.9}
Let $|\cdot|$ denote the standard Euclidean norm in $\mathbb{R}^n$ and let
$f : \mathbb{R}^n \to [0,\infty)$ be another norm (that is, a function satisfying the properties
of \textit{Definition 9.89}).
\begin{enumerate}[label=\textbf{\arabic*.}]
  \item Expressing $x$ in a basis and using the ``abstract'' properties that $f$ must have, show
    that there is a constant $C_1 > 0$ such that $f(x) \leq C_1|x|$ for all $x \in \mathbb{R}^n$.
  \item Show that $f$ is continuous in $\mathbb{R}^n$ (with respect to the standard distance of
    $\mathbb{R}^n$!).
  \item Show that there is a number $c_2 > 0$ such that $f(x) \geq c_2$ for all $|x| = 1$.
  \item Conclude that $f(x) \geq c_2|x|$ for all $x \in \mathbb{R}^n$.
  \item Show that if $\tilde f$ is yet another norm, then there is $C > 0$ such that
    $C^{-1}f(x) \leq \tilde f(x) \leq Cf(x)$ for all $x \in \mathbb{R}^n$.
\end{enumerate}
\exhint[Official hints]{\textbf{3.9.2:} recall that from the definition it follows that
$|f(x) - f(y)| \leq f(x-y)$, and that Lipschitz functions are always continuous.
\textbf{3.9.3:} apply the Weierstrass theorem to $f$ on $S := \{x \in \mathbb{R}^n : |x| = 1\}$
and use that $f$, being a norm, is nondegenerate. \textbf{3.9.5:} if $f$ is equivalent to $|\cdot|$ and
$\tilde f$ is equivalent to $|\cdot|$, it follows by transitivity that $f$ is equivalent to
$\tilde f$.}
\end{exercise}

\begin{remark}
Hint 3.9.3 is exactly Corsin's reason \#2 for caring about compactness
(\cref{sec:why_compactness_matters}), namely the Weierstrass extreme value theorem applied on the
compact unit sphere.
\end{remark}

% Supplement: Analysis_II_Script_v1.pdf, p. 35
\begin{exercise}[The polynomial ring $\mathbb{R}{[}x{]}$ is infinite-dimensional]
\label{ex:polynomials_infinite_dimensional}
Show that the vector space $\mathbb{R}[x]$ of polynomials with real coefficients is infinite
dimensional, and exhibit a basis.
\end{exercise}

\begin{remark}
This is exactly the hypothesis \cref{ex:3.9} needs and cannot do without.
\Cref{rem:norm_equivalence_needs_finite_dimension} already gave $C^0([0,1],\mathbb{R})$ as an
infinite-dimensional space where norm equivalence fails. Note that $\mathbb{R}[x]$ is the
algebraically simplest such example, with an explicit basis and no analysis required at all.
\end{remark}

% Originally: content/exercise-sheets/week-03.tex, problem 3.10
% Source: exercises/Ex3_Analysis2_eng.pdf

\begin{exercise}[3.10 --- Hilbert--Schmidt norm of the composition]
\label{ex:3.10}
Take two linear functions $\varphi : \mathbb{R}^d \to \mathbb{R}^n$ and
$\psi : \mathbb{R}^n \to \mathbb{R}^m$, and denote with $\Phi$, $\Psi$ the matrices that represent
them in the canonical bases. Recall that the linear map
$\psi \circ \varphi : \mathbb{R}^d \to \mathbb{R}^m$ is represented in these bases by the matrix
$\Psi \cdot \Phi$. Show that
$$\|\Psi \cdot \Phi\| \leq \|\Psi\|\|\Phi\|,$$
where $\|\cdot\|$ is the Hilbert--Schmidt norm of a matrix (see \textit{Definition 9.96} in the
notes).

\exhint[Official hint]{Recall that $\|M\|^2$ is the sum of the squares of the entries of $M$.
Notice that $(\Psi\cdot\Phi)^i_j$ ($i$-th row, $j$-th column) is the scalar product of $\Psi^i$
and $\Phi_j$, which are vectors of $\mathbb{R}^n$. Apply Cauchy--Schwarz to each of them.}
\end{exercise}

% Originally: content/week-04.tex, \exercisesheet{4}
% Source: exercises/Ex4_Analysis2_eng.pdf

\begin{exercise}[4.3 --- $p$-norms \textnormal{(semi-important; parts 4 and 6 important)}]
\label{ex:4.3}
For $p \geq 1$ and $x \in \mathbb{R}^n$ define the \textbf{$p$-norm} of $x$ as
\[|x|_p := \Big(\sum_{i=1}^{n} |x_i|^p\Big)^{1/p}.\]
\begin{enumerate}[label=\textbf{\arabic*.}]
  \item For $n = 2$ and $p = 1, 2, 10$ sketch the sets $\{x \in \mathbb{R}^2 : |x|_p \leq 1\}$.
  \item For a given $x \in \mathbb{R}^n$, compute the limit
    $|x|_\infty := \lim_{p\to\infty}|x|_p$.
  \item Using an appropriate inequality that you have seen in class, prove that
    \[\Big(\sum_{i=1}^{n} a_i^{p-1}b_i\Big)^2 \leq \Big(\sum_{i=1}^{n}a_i^p\Big)\Big(\sum_{i=1}^{n}a_i^{p-2}b_i^2\Big),\]
    whenever $a_i, b_i$ are $n$-tuples of positive numbers.
  \item Fix $x, y \in \mathbb{R}^n$ and consider the function $f : [0,1] \to \mathbb{R}$ defined
    as
    \begin{equation*}
    f(t) := |tx + (1-t)y|_p = \Big(\sum_{i=1}^{n}|tx_i + (1-t)y_i|^p\Big)^{1/p}. \tag{1}
    \end{equation*}
    Show that $f$ is convex. You may assume that the coordinates of $x$ and $y$ are all strictly
    positive and use the inequality of the previous point.
  \item Deduce from the previous point that the triangular inequality holds, i.e.\
    $|x+y|_p \leq |x|_p + |y|_p$ for all $x,y \in \mathbb{R}^n$.
  \item What happens for $p \in (0,1)$?
\end{enumerate}
\exhint[Official hints]{\textbf{4.3.3:} use Cauchy--Schwarz and the fact that
$a_i^{p-1}b_i = a_i^{p/2}\cdot a_i^{(p-2)/2}b_i$. \textbf{4.3.4:} show that $f''(t) \geq 0$; it
is not the simplest derivative, but if you get it right the inequality in 4.3.3 will be just what
you need. \textbf{4.3.5:} write the convexity inequality between $x$, $y$ and the middle point
$(x+y)/2$. To get the general case from the one with positive coordinates, observe that for
$x \in \mathbb{R}^n$, denoting $\hat x := (|x_1|,\dots,|x_n|)$, we have
$|x+y|_p \leq |\hat x + \hat y|_p$ and $|\hat x|_p = |x|_p$ for all $x,y \in \mathbb{R}^n$.}
\end{exercise}

% Originally: content/exercise-sheets/week-04.tex, problem 4.4 --- never previously built
% Source: exercises/Ex4_Analysis2_eng.pdf

\begin{exercise}[4.4 --- $p$-means \textnormal{(optional)}]
\label{ex:4.4}
For $x \in \mathbb{R}^n$ with positive coordinates and $p \neq 0$ define the \textbf{$p$-mean} as
$$\mu_p(x) := \left(\frac{x_1^p + \dots + x_n^p}{n}\right)^{1/p}.$$
\begin{enumerate}[label=\textbf{\arabic*.}]
  \item Compute the limits $p \to \pm\infty$, $p \to 0$ and define accordingly
    $\mu_{-\infty}(x)$, $\mu_0(x)$, $\mu_{+\infty}(x)$.
  \item For any $n$-tuple of numbers $a_i > 0$ show that
    $$\sum_{i=1}^{n} \frac{a_i\log(a_i)}{a_1 + \dots + a_n} \geq \log\left(\frac{a_1 \dots + a_n}{n}\right).$$
  \item For a fixed $x$, show that the function $f : \mathbb{R}\to\mathbb{R}$, given by
    $f(t) := \mu_t(x)$, is continuous and increasing.
  \item Prove the Arithmetic--Geometric inequality and Arithmetic--Quadratic inequality:
    $$n(x_1x_2\cdots x_n)^{1/n} \leq x_1 + \dots + x_n, \qquad (x_1+\dots+x_n)^2 \leq n(x_1^2+\dots+x_n^2).$$
  \item \textbf{(*)} Is $f$ continuously differentiable in the whole $\mathbb{R}$?
\end{enumerate}
\exhint[Official hints]{\textbf{4.4.2:} combine the concavity of $\log(\cdot)$ and the
Cauchy--Schwarz inequality; use the generalisation of the concavity inequality: for any concave
$f : \mathbb{R}\to\mathbb{R}$,
$f(\lambda_1x_1 + \dots + \lambda_nx_n) \geq \lambda_1f(x_1)+\dots+\lambda_nf(x_n)$ for all
$x_i \in \mathbb{R}$, $0 \leq \lambda_i \leq 1$ with $\lambda_1+\dots+\lambda_n = 1$.
\textbf{4.4.3:} it might be more convenient to work with $\log f(t)$.}
\end{exercise}
